\section{统一分数拒绝虚假加速}

$Q_{\mathrm{HW}}$ 将性能增益和资源代价聚合为单一有界效用。令锚点 $a$ 和
候选方案 $c$ 分别具有周期延迟 $L$、有效周期 $p$ 和启动间隔 $II$。性能比
$r_P = (p_a L_a)/(p_c L_c)$ 捕获加速效果；面积比 $r_A$ 使用容量归一化占用
$F(x)=\sum_{r} x_r/K_r$（覆盖 LUT、FF、DSP、BRAM、URAM）防止资源转移导致
分数跳变。两项证据因子奖励源码修改（$D$，$1.01{\times}$）和有效性修复
（$F$，$2{\times}$）。综合硬件比
$h = 1.01^{D}\,2^{F}\,r_P^{\,0.55}\,r_A^{\,0.45}$ 以 55\% 权重偏向性能、
45\% 偏向面积，基于 99 个冻结 PPA 参考校准。

$Q_{\mathrm{HW}} = 1-1/(1+h)^2$ 将 $h{=}1$（无变化）映射为 $0.75$，
$h{\to}\infty$ 映射为 $1$，$h{\to}0$ 映射为 $0$，对膨胀加速连续衰减。预算
效率 $E = \max(0.80,\,1-0.10u_{\mathrm{credit}}-0.10u_{\mathrm{time}})$ 对
过多工具调用折价，保底 $0.80$，其中各 $u$ 为消耗比例截断至 $[0,1]$。最终
分数为
\begin{equation}
S = 100\,V(c)\,Q_{\mathrm{HW}}\,E.
\label{eq:score}
\end{equation}
提升仅在 $V(c){=}1$ 后比较 $Q_{\mathrm{HW}}$；门失败得 $S{=}0$，效率保底
确保有效基线至少获得 $60$ 分。
